\chapter{Interrupt Controller}
\label{chap:interrupts}

The 6809 has three interrupt inputs: IRQ, FIRQ and NMI. On the \jrtwo\ and \ktwo\ nearly every device interrupts through IRQ. FIRQ isn't connected, and NMI comes only from the IEC port, when you ask for it (section~\ref{sec:int_nmi}).

Many devices share the one IRQ line, so an interrupt controller sits between them and the CPU. It records which devices have raised an interrupt, lets you choose which of them reach the CPU, and lets your handler find the cause without polling every device.

If you write programs for NitrOS-9, the system and its drivers own the interrupt controller; read section~\ref{sec:int_nos9}. If you write raw 6809 programs, you set it up yourself.

\section{Groups and registers}

The controller handles up to 32 sources in four groups of eight. Each group has four registers, and each source has the same bit position in all four:

\begin{description}
    \item[PENDING] reads 1 for each source that has raised an interrupt. Write 1 to a bit to clear it; writing 0 leaves it alone, so you can clear one source without touching the others.
    \item[MASK] decides which pending flags interrupt the CPU. A masked source (1) still sets its pending flag, but doesn't interrupt; an unmasked source (0) does.
    \item[EDGE] and \textbf{POLARITY} decide what sets the pending flag (table~\ref{tab:int_pol_edge}).
\end{description}

The controller holds IRQ low while any pending flag is set and unmasked. Your handler must clear the pending flags it services, or the CPU takes the interrupt again as soon as it returns.

Table~\ref{tab:int_reg} shows the registers and the source of each bit, and the list after it describes them. The \ktwo\ has three sources the \jrtwo\ lacks: VIA 1 for the built-in keyboard, the optical keyboard and the WizNet FIFO. On the \jrtwo\ those bits never become pending.

% Generated by tools/regs.py from data/registers/interrupts.toml. Do not edit.
\begin{table}[ht]
    \begin{center}
        \setlength{\tabcolsep}{3pt}
        \begin{tabular}{|c|l|BBBBBBBB} \hline
            Address & Name & 7 & 6 & 5 & 4 & 3 & 2 & 1 & 0 \\ \hline\hline
            \verb+$FE20+ & INT\_\allowbreak{}PENDING\_\allowbreak{}0 & \bitspan{1}{INT\_\allowbreak{}CARTRIDGE} & \bitspan{1}{INT\_\allowbreak{}DMA} & \bitspan{1}{INT\_\allowbreak{}TIMER\_\allowbreak{}1} & \bitspan{1}{INT\_\allowbreak{}TIMER\_\allowbreak{}0} & \bitspan{1}{INT\_\allowbreak{}PS2\_\allowbreak{}MOUSE} & \bitspan{1}{INT\_\allowbreak{}PS2\_\allowbreak{}KBD} & \bitspan{1}{INT\_\allowbreak{}VKY\_\allowbreak{}SOL} & \bitspan{1}{INT\_\allowbreak{}VKY\_\allowbreak{}SOF} \\ \hline
            \verb+$FE21+ & INT\_\allowbreak{}PENDING\_\allowbreak{}1 & \bitspan{1}{INT\_\allowbreak{}SDC\_\allowbreak{}INS} & \bitspan{1}{INT\_\allowbreak{}VIA1} & \bitspan{1}{INT\_\allowbreak{}VIA0} & \bitspan{1}{INT\_\allowbreak{}RTC} & \bitspan{3}{---} & \bitspan{1}{INT\_\allowbreak{}UART} \\ \hline
            \verb+$FE22+ & INT\_\allowbreak{}PENDING\_\allowbreak{}2 & \bitspan{1}{---} & \bitspan{1}{INT\_\allowbreak{}HDMI\_\allowbreak{}PIN} & \bitspan{1}{INT\_\allowbreak{}WIFI\_\allowbreak{}PIN} & \bitspan{1}{INT\_\allowbreak{}NET\_\allowbreak{}PIN} & \bitspan{1}{IEC\_\allowbreak{}SREQ\_\allowbreak{}i} & \bitspan{1}{IEC\_\allowbreak{}ATN\_\allowbreak{}i} & \bitspan{1}{IEC\_\allowbreak{}CLK\_\allowbreak{}i} & \bitspan{1}{IEC\_\allowbreak{}DATA\_\allowbreak{}i} \\ \hline
            \verb+$FE23+ & INT\_\allowbreak{}PENDING\_\allowbreak{}3 & \bitspan{2}{---} & \bitspan{1}{INT\_\allowbreak{}WIZFI\_\allowbreak{}TX} & \bitspan{1}{---} & \bitspan{1}{INT\_\allowbreak{}WIZNET} & \bitspan{1}{INT\_\allowbreak{}OPT\_\allowbreak{}KBD} & \bitspan{1}{INT\_\allowbreak{}MIDI\_\allowbreak{}RX} & \bitspan{1}{INT\_\allowbreak{}WIZFI\_\allowbreak{}RX} \\ \hline
            \verb+$FE24+--\verb+$FE27+ & INT\_\allowbreak{}POLARITY\_\allowbreak{}0–3 & \bitspan{8}{One bit per source} \\ \hline
            \verb+$FE28+--\verb+$FE2B+ & INT\_\allowbreak{}EDGE\_\allowbreak{}0–3 & \bitspan{8}{One bit per source} \\ \hline
            \verb+$FE2C+--\verb+$FE2F+ & INT\_\allowbreak{}MASK\_\allowbreak{}0–3 & \bitspan{8}{One bit per source} \\ \hline
        \end{tabular}
    \end{center}
    \caption{Interrupt controller registers}
    \label{tab:int_reg}
\end{table}

% Generated by tools/regs.py from data/registers/interrupts.toml. Do not edit.
\begin{description}
    \item[INT\_PENDING\_0] (\verb+$FE20+, read/write, reset \verb+$00+) Group 0 pending flags. A bit reads 1 when its source has raised an interrupt. Write 1 to a bit to clear it; writing 0 leaves it alone.
    \begin{description}
        \item[INT\_CARTRIDGE] (bit 7) Interrupt line on the cartridge connector.
        \item[INT\_DMA] (bit 6) DMA transfer finished (when enabled in the DMA control register).
        \item[INT\_TIMER\_1] (bit 5) Timer 1 reached its target value.
        \item[INT\_TIMER\_0] (bit 4) Timer 0 reached its target value.
        \item[INT\_PS2\_MOUSE] (bit 3) PS/2 mouse data.
        \item[INT\_PS2\_KBD] (bit 2) PS/2 keyboard data.
        \item[INT\_VKY\_SOL] (bit 1) VICKY start of line: the display reached the line set in the line interrupt registers.
        \item[INT\_VKY\_SOF] (bit 0) VICKY start of frame.
    \end{description}
    \item[INT\_PENDING\_1] (\verb+$FE21+, read/write, reset \verb+$00+) Group 1 pending flags. Bits 3–1 have no source.
    \begin{description}
        \item[INT\_SDC\_INS] (bit 7) An SD card was inserted.
        \item[INT\_VIA1] (bit 6) VIA 1, which scans the K2's built-in keyboard. Not connected on the Jr2.
        \item[INT\_VIA0] (bit 5) VIA 0 (the joystick ports).
        \item[INT\_RTC] (bit 4) Real-time clock.
        \item[INT\_UART] (bit 0) The serial port (UART).
    \end{description}
    \item[INT\_PENDING\_2] (\verb+$FE22+, read/write, reset \verb+$00+) Group 2 pending flags: the IEC bus lines and the interrupt pins of the network, Wi-Fi and HDMI parts. Bit 7 has no source. The IEC bits can also raise an NMI (see the text).
    \begin{description}
        \item[INT\_HDMI\_PIN] (bit 6) HDMI interface interrupt pin.
        \item[INT\_WIFI\_PIN] (bit 5) Wi-Fi module interrupt pin.
        \item[INT\_NET\_PIN] (bit 4) Network module interrupt pin.
        \item[IEC\_SREQ\_i] (bit 3) IEC SREQ line.
        \item[IEC\_ATN\_i] (bit 2) IEC ATN line.
        \item[IEC\_CLK\_i] (bit 1) IEC clock line.
        \item[IEC\_DATA\_i] (bit 0) IEC data line.
    \end{description}
    \item[INT\_PENDING\_3] (\verb+$FE23+, read/write, reset \verb+$00+) Group 3 pending flags: receive and transmit FIFOs. Bit 4 (INT\_MIDI\_VS\_RX in NitrOS-9) and bits 7–6 have no source on either board.
    \begin{description}
        \item[INT\_WIZFI\_TX] (bit 5) The Wi-Fi transmit FIFO became empty.
        \item[INT\_WIZNET] (bit 3) The WizNet FIFO has data. K2 only.
        \item[INT\_OPT\_KBD] (bit 2) The optical keyboard FIFO has data. K2 only.
        \item[INT\_MIDI\_RX] (bit 1) The MIDI receive FIFO has data.
        \item[INT\_WIZFI\_RX] (bit 0) The Wi-Fi receive FIFO has data.
    \end{description}
    \item[INT\_POLARITY\_0–3] (\verb+$FE24+--\verb+$FE27+, read/write, reset \verb+$00+) Polarity for groups 0–3 (\$FE24–\$FE27), one bit per source in the same positions as the pending flags. With EDGE, chooses the condition that sets the pending flag.
    \item[INT\_EDGE\_0–3] (\verb+$FE28+--\verb+$FE2B+, read/write, reset \verb+$FF+) Edge select for groups 0–3 (\$FE28–\$FE2B). 1 = the flag is set by a change of the input; 0 = it is set while the input is at a level.
    \item[INT\_MASK\_0–3] (\verb+$FE2C+--\verb+$FE2F+, read/write, reset \verb+$FF+) Masks for groups 0–3 (\$FE2C–\$FE2F). 1 = the source can set its pending flag but doesn't interrupt the CPU; 0 = a pending flag interrupts the CPU.
\end{description}


\begin{note}
    Some devices have an interrupt enable of their own as well, such as the VIAs, the DMA engine and the line interrupt. Such a device interrupts the CPU only when you enable it in the device and unmask it in the controller.
\end{note}

\subsection*{Edge and polarity}

\begin{table}[ht]
    \begin{center}
        \begin{tabular}{|c|c|l|} \hline
            EDGE & POLARITY & The pending flag is set \\ \hline\hline
            0 & 0 & while the input is low \\ \hline
            0 & 1 & while the input is high \\ \hline
            1 & 0 & when the input goes from high to low \\ \hline
            1 & 1 & when the input goes from low to high \\ \hline
        \end{tabular}
    \end{center}
    \caption{Interrupt polarity and edge}
    \label{tab:int_pol_edge}
\end{table}

After reset every source is edge-triggered on a falling edge (EDGE \verb+$FF+, POLARITY \verb+$00+). That suits all the device sources, and NitrOS-9 never changes it. With a level setting, the flag is set again as long as the input stays at that level, even straight after you clear it. The edge and polarity registers matter mainly for the IEC lines in group 2, where you may want to catch a particular change on the bus.

\subsection*{Start of frame and start of line}

VICKY raises the start-of-frame interrupt (SOF) once per frame, when it has drawn the last visible line, which makes it the right moment to move sprites or change colors without tearing. It also makes a steady clock: \verify{60 times a second in the 640×480 modes}.

The start-of-line interrupt (SOL) happens when the display reaches a line you choose. Set the line and enable the interrupt in VICKY's line interrupt registers (table~\ref{tab:lint_reg}); then unmask INT\_VKY\_SOL.

\section{Reset state}

At power-up and reset the controller masks every source, clears every pending flag, and sets every source to falling-edge triggering. The 6809 itself starts with IRQ and FIRQ masked in its condition code register.

Unmask a source only after you have installed its handler and cleared its pending flag; a flag left over from before can otherwise interrupt at once.

\section{NMI from the IEC port}
\label{sec:int_nmi}

Bit 5 of the IEC output register (\verb+$FE81+) routes the IEC interrupts to NMI as well. While the bit is set, any unmasked pending IEC flag (group 2, bits 0--3) pulls NMI low as well as IRQ. While it is clear, the IEC flags use IRQ only. No other source can raise NMI. \verify{The bit is clear after reset.}

The 6809 ignores NMI until your program has loaded S for the first time. NMI can't be masked in the CPU, so mask the IEC sources in INT\_MASK\_2 when you don't want it.

\section{Raw 6809 example: the start-of-frame interrupt}

This example counts frames in the top left character of the screen. It builds on the raw skeleton in \verb+examples/raw/skeleton.asm+, which already masks and clears every interrupt and puts the 6809 vectors at the end of its image. It uses the names in \verb+examples/include/wildbits_regs.inc+, plus \verb+INT_VKY_SOF equ $01+.

The handler maps the text screen (block \verb+$FE+) into slot 5 to change the character, then puts slot 5 back. The 6809 saves every register when it takes an IRQ, so the handler needs to save only the slot it changes.

\begin{verbatim}
IrqHandler  lda   #INT_VKY_SOF
            bita  INT_PENDING_0     is SOF pending?
            beq   @done
            sta   INT_PENDING_0     yes: clear it
            lda   MMU_SLOT_5        save slot 5 ($A000-$BFFF)
            pshs  a
            lda   #$FE              text screen block
            sta   MMU_SLOT_5
            inc   $A000             the character at row 0, column 0
            puls  a
            sta   MMU_SLOT_5        put slot 5 back
@done       rti
\end{verbatim}

Point the IRQ vector at the handler in the skeleton's vector table (\verb+fdb IrqHandler+ at \verb+$FFF8+), and turn the interrupt on in your setup code:

\begin{verbatim}
            orcc  #$10              hold off IRQs while we set up
            lda   #INT_VKY_SOF
            sta   INT_PENDING_0     clear any old SOF
            lda   INT_MASK_0
            anda  #~INT_VKY_SOF     unmask SOF only
            sta   INT_MASK_0
            andcc #~$10             let IRQs in
\end{verbatim}

The handler reads \verb+MMU_SLOT_5+ to save it; that read returns the LUT selected by EDIT\_LUT, so this example assumes EDIT\_LUT and ACT\_LUT are the same, as after reset.

\section{Under NitrOS-9}
\label{sec:int_nos9}

When NitrOS-9 starts, it masks every source and clears every pending flag. Each driver then unmasks the sources it handles, changing only its own bits with a read-modify-write.

Interrupts belong to device drivers. A driver installs its handler with \verb+F$IRQ+, giving the address of the pending register (\verb+INT_PENDING_0+ to \verb+INT_PENDING_3+) and a packet with a flip byte, a mask byte for its source, and a priority. When an IRQ arrives, NitrOS-9 checks each installed handler's pending register against its mask and calls the handlers whose bits are set. The handler clears its pending flag and returns. This is how the start-of-line driver installs its handler:

\begin{verbatim}
            ldd   #INT_PENDING_0    the register to check
            leax  SOL_Pckt,pcr      flip, mask and priority
            leay  SOL_IRQSvc,pcr    the handler
            os9   F$IRQ
            ...
            lda   INT_MASK_0        unmask SOL, keeping the other bits
            anda  #^INT_VKY_SOL
            sta   INT_MASK_0

SOL_Pckt    fcb   %00000000         flip byte
            fcb   INT_VKY_SOL       mask byte
            fcb   $F1               priority
\end{verbatim}

A program doesn't handle interrupts itself. It asks a driver to send it a signal when the event happens, and handles the signal with an intercept routine (\verb+F$Icpt+). For example, the \verb+SS.SOLIRQ+ SetStat asks the start-of-line driver to signal your process when the display reaches a given line. Don't write the interrupt controller's registers from a program: you would take sources away from the drivers that own them.
